\section{Preliminaries: System Model and Problem Statement}
\label{sec:system}

This section defines the downlink frame, the channel and the reports the base station receives,
the progressive code, the digital chain it is compared with, and the scheduling problem.
Table~\ref{tab:notation} collects the notation, with the default values of
Section~\ref{sec:eval}, and Fig.~\ref{fig:system} illustrates the frame and the code.

\begin{table}[!t]
\caption{Notation. Values are the defaults of Section~\ref{sec:eval};
$\mathcal{U}[a, b]$ is the uniform distribution on $[a, b]$.}
\label{tab:notation}
\centering
\footnotesize
\setlength{\tabcolsep}{4pt}
\renewcommand{\arraystretch}{1.1}
\begin{tabular}{|c|>{\raggedright\arraybackslash}p{4.45cm}|>{\raggedright\arraybackslash}p{1.75cm}|}
\hline
Symbol & Meaning & Value \\
\hline
$K$, $\mathcal{K}$, $k$ & number, set and index of users & $K = 8$ \\
$T$, $t$ & deadline in slots, slot index & $T = 24$ \\
$B$ & resource units per slot & $1$ \\
$B_t$ & units available in slot $t$ & \\
$p$ & probability that a unit is taken & $0$ \\
$S$ & complex channel uses per unit & $3HW/48$ \\
$a_{k,t}$ & $1$ if user $k$ is served in slot $t$ & \\
$\mathbf{a}_t$ & schedule $(a_{1,t}, \dots, a_{K,t})$ of slot $t$ & \\
\hline
$\bar\gamma_k$ & mean SNR of user $k$ & $\mathcal{U}[0, 15]$~dB \\
$h_{k,t}$ & fading coefficient & $\mathcal{CN}(0,1)$ \\
$\gamma_{k,t}$, $n_{k,t}$ & slot SNR, noise variance & \\
$\gamma_{\min}$, $\gamma_{\max}$ & SNR range of training & $-2$, $22$~dB \\
$\hat\gamma_{k,t}$, $\hat h_{k,t}$ & reported SNR, reported fading & \\
$\epsilon_{k,t}$, $\sigma_\epsilon$ & report error, standard deviation & $\sigma_\epsilon = 0$--$3$~dB \\
$r$, $f_D$, $\tau$ & correlation, Doppler frequency, age & $r = 1$--$0$ \\
\hline
$\mathbf{x}_k$, $\hat{\mathbf{x}}_k$ & image, its reconstruction & \\
$H$, $W$ & image height and width in pixels & \\
$f$, $g$ & encoder, decoder & \\
$P$, $j$ & phases per code, phase index & $P = 6$ \\
$\mathbf{u}_{k,j}$, $\mathbf{z}_{k,j}$ & phase $j$, as sent and as received & \\
$t_{k,j}$ & slot that carries phase $j$ & \\
$L_k$ & phases held at the deadline & $L_k \le P$ \\
$\gamma_{k,(j)}$, $n_{k,(j)}$ & SNR, noise variance of phase $j$ & \\
$D_k$, $q_k$ & MSE, PSNR at the deadline & \\
$D_\emptyset$ & MSE with no phase received & $0.1$ \\
\hline
$c$, $\eta_c$ & CQI index, its efficiency & \cite{TS36213} \\
$\theta_c$ & SNR threshold of index $c$ & $-6.7$--$22.7$~dB \\
$R(\cdot)$ & rate selected for an SNR & \\
$\eta_{k,t}$ & rate of the block in slot $t$ & \\
$\beta$ & rate backoff & tuned \\
$M$ & transmissions per block, HARQ & $4$ \\
$I(\cdot)$ & information per transmission & \\
$\Gamma$ & SNR gap of the idealized code & $3$~dB \\
$b_k$ & bits per pixel received & \\
\hline
$\pi$, $\mathcal{I}_t$ & policy, information at slot $t$ & \\
$Q$ & quality target & $30$~dB \\
$J_{\mathrm{mean}}$, $J_Q$ & mean PSNR, share at the target & \\
\hline
\end{tabular}
\end{table}

\subsection{Downlink Frame}
\label{sec:system-frame}

A base station serves $K$ users, indexed by $k \in \mathcal{K} = \{1, \dots, K\}$, over a frame of
$T$ slots, indexed by $t = 1, \dots, T$. Each user is owed one image
$\mathbf{x}_k \in [0,1]^{3 \times H \times W}$, of $H \times W$ pixels with three color channels and
intensities scaled to $[0,1]$, and every image must be delivered by the end of slot $T$, the
deadline. Each slot holds $B$ resource units of $S$ complex channel uses each: $B = 1$ is
time-division multiple access, and $B > 1$ models $B$ groups of subbands within one slot. A user
receives at most one unit per slot. Other traffic may take a unit: each unit is available to the
frame with probability $1 - p$, independently of the other units and slots, so the number of
units available in slot $t$ is
\begin{equation}
B_t \sim \mathrm{Binomial}(B, 1 - p),
\label{eq:units}
\end{equation}
and $B_t$ becomes known at the start of slot $t$. The schedule is a set of binary variables
$a_{k,t} \in \{0, 1\}$, with $a_{k,t} = 1$ if user $k$ receives a unit in slot $t$. Since a user
takes at most one unit, a slot can serve at most $B_t$ users:
\begin{equation}
\sum_{k=1}^{K} a_{k,t} \le B_t, \quad t = 1, \dots, T.
\label{eq:capacity}
\end{equation}

\subsection{Channel and Channel Reports}
\label{sec:system-channel}

The channel of each user is flat within a slot and changes independently from slot to slot
(block fading). In slot $t$, a complex symbol $s$ of unit average power sent to user $k$ arrives
as
\begin{equation}
y = \sqrt{\bar\gamma_k}\, h_{k,t}\, s + w,
\label{eq:channel}
\end{equation}
where $\bar\gamma_k$ is the mean SNR of user $k$, set by its path loss;
$h_{k,t} \sim \mathcal{CN}(0, 1)$ is its Rayleigh fading coefficient, independent over users and
slots; and $w \sim \mathcal{CN}(0, 1)$ is the receiver noise, independent over symbols. $\mathcal{CN}(\mu, v)$ denotes a circularly symmetric complex Gaussian law of mean
$\mu$ and variance $v$. The receiver knows $h_{k,t}$ and equalizes,
$z = y / (\sqrt{\bar\gamma_k}\, h_{k,t}) = s + w / (\sqrt{\bar\gamma_k}\, h_{k,t})$, so every
symbol of the slot crosses an additive white Gaussian noise (AWGN) channel with SNR
$\gamma_{k,t}$ and noise variance $n_{k,t}$:
\begin{equation}
\gamma_{k,t} = \bar\gamma_k\, |h_{k,t}|^2, \qquad n_{k,t} = \frac{1}{\gamma_{k,t}}.
\label{eq:slot-snr}
\end{equation}
As $|h_{k,t}|^2$ is exponential with unit mean, $\gamma_{k,t}$ is exponential with mean
$\bar\gamma_k$. A superscript dB marks a value in decibels, as in
$\gamma^{\mathrm{dB}}_{k,t} = 10 \log_{10} \gamma_{k,t} = \bar\gamma^{\mathrm{dB}}_k +
10 \log_{10} |h_{k,t}|^2$.

At the start of slot $t$, the base station holds a report $\hat\gamma_{k,t}$ of the slot SNR of
every user, and nothing about later slots. Three report models are considered.
\emph{Exact reports} give $\hat\gamma_{k,t} = \gamma_{k,t}$. \emph{Erroneous reports} are off by
a Gaussian error in decibels,
\begin{equation}
\hat\gamma^{\mathrm{dB}}_{k,t} = \gamma^{\mathrm{dB}}_{k,t} + \epsilon_{k,t}, \qquad
\epsilon_{k,t} \sim \mathcal{N}(0, \sigma_\epsilon^2),
\label{eq:cqi-error}
\end{equation}
independent over users and slots, where $\mathcal{N}(\mu, v)$ denotes a real Gaussian law of mean
$\mu$ and variance $v$, so that $\sigma_\epsilon$ is the standard deviation of the error.
\emph{Aged reports} were measured a time $\tau$ before the slot, and the channel has changed
since:
\begin{equation}
\hat h_{k,t} = r\, h_{k,t} + \sqrt{1 - r^2}\, e_{k,t}, \qquad
\hat\gamma_{k,t} = \bar\gamma_k\, |\hat h_{k,t}|^2,
\label{eq:aged}
\end{equation}
where $\hat h_{k,t}$ is the fading coefficient the report was measured on,
$e_{k,t} \sim \mathcal{CN}(0, 1)$ is independent of $h_{k,t}$, and
$r = |J_0(2\pi f_D \tau)| \in [0, 1]$ is the magnitude of the correlation of the fading process
at lag $\tau$ under Jakes' isotropic-scattering model \cite{Jakes94}, with $f_D$ the Doppler
frequency and $J_0$ the Bessel function of the first kind and order zero. The magnitude suffices:
as $e_{k,t}$ is circularly symmetric, the joint law of $|h_{k,t}|$ and $|\hat h_{k,t}|$, which is
all the SNRs depend on, is the same for a correlation of either sign. A report with $r = 1$ is
fresh, and one with $r = 0$ says nothing about the slot. Since $\hat h_{k,t}$ has unit
variance and $\mathbb{E}[h_{k,t} \hat h_{k,t}^*] = r$, where ${}^*$ denotes complex conjugation,
the two coefficients are jointly Gaussian with correlation $r$, and the channel given the report is
\begin{equation}
h_{k,t} \mid \hat h_{k,t} \sim \mathcal{CN}\big(r\, \hat h_{k,t},\, 1 - r^2\big),
\label{eq:posterior}
\end{equation}
which a transmitter that knows $r$ can use instead of the report itself
(Section~\ref{sec:method-csi}). In the evaluation, slot SNRs and reports are clipped, in
decibels, to the range $[\gamma_{\min}, \gamma_{\max}]$ on which the code was trained.

\begin{figure}[!t]
\centering
\begin{tikzpicture}[x=0.85cm, y=0.55cm, font=\footnotesize, >=stealth,
  u1/.style={fill=blue!18}, u2/.style={fill=orange!28}, u3/.style={fill=green!20}]
\foreach \s/\r/\k/\j in {1/1/2/1, 2/1/3/2, 3/1/2/2, 4/1/3/3, 5/1/2/3, 6/1/2/4,
                         1/0/3/1, 2/0/1/1, 3/0/1/2, 5/0/1/3, 6/0/3/4} {
  \filldraw[u\k, draw=black] (\s-1, \r) rectangle (\s, \r+1);
  \node at (\s-0.5, \r+0.5) {$\mathbf{u}_{\k,\j}$};
}
\fill[pattern=north east lines, pattern color=black!50] (3, 0) rectangle (4, 1);
\draw (3, 0) rectangle (4, 1);
\foreach \s in {2, 3, 5} \draw[line width=1.1pt] (\s-1, 0) rectangle (\s, 1);
\foreach \s in {1, ..., 6} \node at (\s-0.5, 2.5) {$\s$};
\node[anchor=east] at (-0.15, 2.5) {$t$};
\draw[decorate, decoration={brace, amplitude=3pt}] (-0.15, 0) -- (-0.15, 2)
  node[midway, left=3pt] {$B$};
\draw[dashed] (6, -2.1) -- (6, 3.0) node[above] {deadline};
\node[draw, inner sep=2.5pt] (g) at (7.0, -0.85) {$g$};
\node (xh) at (8.1, -0.85) {$\hat{\mathbf{x}}_1$};
\draw[->] (0, -0.85) -- (g);
\draw[->] (g) -- (xh);
\node[anchor=east] at (-0.15, -0.85) {user 1};
\foreach \s/\j in {2/1, 3/2, 5/3} {
  \node[fill=white, inner sep=1.5pt] (z\j) at (\s-0.5, -0.85) {$\mathbf{z}_{1,\j}$};
  \draw[->] (\s-0.5, 0) -- (z\j);
  \node at (\s-0.5, -1.65) {$\gamma_{1,(\j)}$};
}
\end{tikzpicture}
\caption{A frame with $K = 3$ users, $B = 2$ units per slot and $T = 6$ slots. Each cell is a
resource unit, shaded by the user it serves and labeled with the phase it carries; the column of
slot $t$ is the schedule $\mathbf{a}_t$, chosen at the start of the slot from $\mathcal{I}_t$.
Other traffic takes the hatched unit, so $B_4 = 1$. User~1 is served in slots 2, 3 and 5, so
$t_{1,1} = 2$, $t_{1,2} = 3$, $t_{1,3} = 5$ and $L_1 = 3$; its phases arrive as
$\mathbf{z}_{1,j}$, each at the SNR $\gamma_{1,(j)} = \gamma_{1, t_{1,j}}$ of its slot, and the
decoder $g$ reconstructs $\hat{\mathbf{x}}_1$ from them at the deadline.}
\label{fig:system}
\end{figure}

\subsection{The Progressive Code}
\label{sec:system-code}

Each image is encoded once, before the frame, by an encoder $f$ that maps it to $P$ phases of $S$
complex symbols each,
\begin{equation}
f(\mathbf{x}_k) = (\mathbf{u}_{k,1}, \dots, \mathbf{u}_{k,P}), \qquad \mathbf{u}_{k,j} \in
\mathbb{C}^S,
\label{eq:encoder}
\end{equation}
normalized to unit average power per complex symbol. A phase fills one resource unit, so it has
a bandwidth ratio, the number of complex channel uses per real source dimension, of
$S / (3HW)$, and the whole code has $P S / (3HW)$; with the values of Table~\ref{tab:notation},
these are $1/48$ and $1/8$. The phases are sent in order: the $j$-th unit that user $k$ receives
carries $\mathbf{u}_{k,j}$. User $k$ receives $L_k = \sum_{t=1}^{T} a_{k,t}$ units in the frame,
at most $P$, so it holds $L_k$ phases at the deadline, and for $j = 1, \dots, L_k$ the slot that
carries phase $j$ is
\begin{equation}
t_{k,j} = \min \Big\{ t : \sum_{t' = 1}^{t} a_{k,t'} = j \Big\},
\label{eq:phase-slot}
\end{equation}
the first slot by which user $k$ has been served $j$ times (Fig.~\ref{fig:system}). Phase $j$
arrives after equalization as
\begin{equation}
\mathbf{z}_{k,j} = \mathbf{u}_{k,j} + \mathbf{w}_{k,j}, \qquad
\mathbf{w}_{k,j} \sim \mathcal{CN}\big(\mathbf{0},\, n_{k,(j)} \mathbf{I}_S\big),
\label{eq:phase-rx}
\end{equation}
where $\mathbf{I}_S$ is the $S \times S$ identity matrix, so the $S$ entries of the noise
$\mathbf{w}_{k,j}$ are independent with variance $n_{k,(j)} = n_{k, t_{k,j}}$, that of the slot
that carried the phase. Each phase thus meets the SNR of its own slot,
$\gamma_{k,(j)} = \gamma_{k, t_{k,j}}$. At the deadline, user $k$ decodes the prefix it holds with
one decoder $g$ that accepts any number of phases,
\begin{equation}
\hat{\mathbf{x}}_k = g(\mathbf{z}_{k,1}, \dots, \mathbf{z}_{k,L_k}),
\label{eq:decoder}
\end{equation}
and neither $f$ nor $g$ is told an SNR. The quality of the reconstruction is its mean squared
error (MSE) per pixel and color channel, $D_k$, or its peak signal-to-noise ratio (PSNR), $q_k$:
\begin{equation}
D_k = \frac{1}{3HW} \big\| \mathbf{x}_k - \hat{\mathbf{x}}_k \big\|_2^2, \qquad
q_k = 10 \log_{10} \frac{1}{D_k}~\text{dB},
\label{eq:psnr}
\end{equation}
where the peak intensity is 1. A user that receives no phase is assigned $D_k = D_\emptyset$,
a PSNR of $10$~dB for the value of Table~\ref{tab:notation}. The decoded quality thus depends on
the image, on the number of phases received and on the SNR each of them met. For a trained,
nonlinear decoder no closed form of this dependence is known;
Section~\ref{sec:method-abstraction} models it.

\subsection{The Digital Baseline}
\label{sec:system-digital}

Physical link design sends the image as bits. The image is compressed into HEIC files, which use
HEVC intra coding \cite{Sullivan12}, at several sizes. A unit given to user $k$ in slot $t$ carries one
block of $S$ channel uses at a spectral efficiency $\eta_{k,t}$, in bits per channel use, set from
the report through the CQI table:
\begin{equation}
\begin{aligned}
R(\gamma^{\mathrm{dB}}) &= \max \big\{ \eta_c : \theta_c \le \gamma^{\mathrm{dB}} \big\}, \\
\eta_{k,t} &= R\big(\hat\gamma^{\mathrm{dB}}_{k,t} - \beta\big).
\end{aligned}
\label{eq:cqi}
\end{equation}
Here $\eta_c$ is the efficiency of CQI index $c = 1, \dots, 15$ in Table~7.2.3-1 of
\cite{TS36213}, from $0.15$ to $5.55$ bits per channel use. A CQI index is defined there as one
whose transport block could be received with an error probability of at most $0.1$, and
$\theta_c$ is the AWGN SNR we take for that, from $-6.7$~dB for $c = 1$ to $22.7$~dB for
$c = 15$. Below $\theta_1$, $R = 0$ and the unit carries nothing. The backoff $\beta$, in
decibels, trades the rate of a block against the risk that a report error makes it fail: with
exact reports $\beta = 0$, and otherwise Section~\ref{sec:eval} tunes it for each scheduler and
objective. A block is decoded if the true SNR supports its rate,
$R(\gamma^{\mathrm{dB}}_{k,t}) \ge \eta_{k,t}$, and is lost otherwise; an acknowledgment tells
the base station which.

With hybrid automatic repeat request (HARQ), a lost block is sent again in each of the user's
next slots until it is decoded or has been sent $M$ times, when it is dropped; a block still
pending at the deadline is lost. After $m$ transmissions, at SNRs $\gamma_1, \dots, \gamma_m$,
the receiver combines them. With incremental redundancy \cite{Caire01} the block is decoded once
\begin{equation}
\sum_{i=1}^{m} I(\gamma_i) \ge \eta,
\label{eq:harq-ir}
\end{equation}
and with Chase combining \cite{Chase85} once
\begin{equation}
I\Big( \sum_{i=1}^{m} \gamma_i \Big) \ge \eta,
\label{eq:harq-cc}
\end{equation}
where $\eta$ is the rate of the block, set at its first transmission, and $I(\gamma)$, the
information one transmission adds, is the efficiency of the CQI table interpolated linearly in
decibels between its entries; below $\theta_1$ it falls with the slope of the first segment, to
zero at $\theta_1 - 3.7$~dB, and above $\theta_{15}$ it stays at $\eta_{15}$. Because $I(\gamma)$
equals $\eta_c$ at $\theta_c$ and increases strictly up to $\theta_{15}$, a single transmission is
decoded exactly when $R(\gamma^{\mathrm{dB}}) \ge \eta$, as without HARQ.

As with the code, a user is served in at most $P$ slots, so that both deliver an image in at most
$PS$ channel uses. At the deadline user $k$ has received
\begin{equation}
b_k = \frac{S}{HW} \sum_{\text{decoded blocks}} \eta
\label{eq:bpp}
\end{equation}
bits per pixel, and it decodes the largest file of $\mathbf{x}_k$ whose size, in bits per pixel,
is at most $b_k$, or nothing ($D_k = D_\emptyset$) if even the smallest file is larger. An
idealized digital code replaces both $R$ and $I$ by $\log_2(1 + \gamma / \Gamma)$, the rate of a
code whose SNR gap to capacity is $\Gamma$, with the same backoff and HARQ.

\subsection{Problem Statement}
\label{sec:system-problem}

A scheduler is a policy $\pi = (\pi_1, \dots, \pi_T)$ that chooses the schedule of slot $t$ from
what the base station knows at its start:
\begin{equation}
\mathbf{a}_t = \pi_t(\mathcal{I}_t),
\label{eq:policy}
\end{equation}
where $\mathbf{a}_t = (a_{1,t}, \dots, a_{K,t})$ is the schedule of slot $t$ and
\begin{equation}
\mathcal{I}_t = \big\{ B_{t'},\, \hat\gamma_{k,t'} : k \in \mathcal{K},\ t' \le t \big\} \cup
\big\{ \mathbf{a}_{t'} : t' < t \big\}
\label{eq:info}
\end{equation}
holds the units available and the reports up to slot $t$, and the schedules before it. From
$\mathcal{I}_t$ the base station knows how many phases each user holds and in which slots they
were sent, but not the SNRs they met. Besides $\mathcal{I}_t$, it knows the images and the mean
SNRs $\bar\gamma_1, \dots, \bar\gamma_K$ from the start of the frame; a transmitter that accounts
for aged reports also knows $r$; and the digital chain knows the acknowledgments of its past
blocks. The policy is causal: it never sees the report of a later slot. Two objectives are
considered. The first is the mean PSNR at the deadline,
\begin{equation}
J_{\mathrm{mean}}(\pi) = \mathbb{E} \Big[ \frac{1}{K} \sum_{k=1}^{K} q_k \Big],
\label{eq:obj-mean}
\end{equation}
and the second is the share of users whose image reaches a quality target $Q$ by the deadline,
\begin{equation}
J_Q(\pi) = \mathbb{E} \Big[ \frac{1}{K} \sum_{k=1}^{K} \mathbf{1}\{ q_k \ge Q \} \Big],
\label{eq:obj-target}
\end{equation}
where $\mathbf{1}\{\cdot\}$ is the indicator function. The expectations are over the fading, the
reports, the available units and the channel noise, and in the evaluation also over the images
and the mean SNRs drawn for each frame. The scheduling problem is
\begin{equation}
\max_{\pi\ \text{causal}} \ J(\pi) \quad \text{s.t.} \quad \eqref{eq:capacity},\ \
a_{k,t} \in \{0, 1\},\ \ L_k \le P\ \ \forall k,
\label{eq:problem}
\end{equation}
with $J = J_{\mathrm{mean}}$ or $J = J_Q$.

Problem~\eqref{eq:problem} is a finite-horizon stochastic control problem with three
difficulties. First, each user's reward is the decoded quality $q_k$ of \eqref{eq:psnr}, for
which a learned code offers no closed form. Second, the state at slot $t$ holds every user's
delivery history, the number of phases it holds and the SNR each met, so it is continuous, its
dimension grows with $K$, and the users are coupled through \eqref{eq:capacity}. Third, the
decisions are made on the reports $\hat\gamma_{k,t}$, not on the SNRs $\gamma_{k,t}$ that decide
the quality. Section~\ref{sec:method} takes them in turn: a link abstraction predicts $q_k$ and
reduces each user's history to two numbers (Section~\ref{sec:method-abstraction}); the problem
then becomes a restless bandit \cite{Whittle88}, which CLASP relaxes into one exact dynamic
program per user (Section~\ref{sec:method-clasp}); and a calibration-aware reward and the
posterior~\eqref{eq:posterior} account for the reports (Section~\ref{sec:method-csi}).
